\section*{Capítulo \ref{ch:g4:air-a-mixture-of-gases} --- El aire, una mezcla de gases}

\begin{solution}{exo:g4:air-a-mixture-of-gases:1}
Una mezcla. Sus dos gases principales son el nitrógeno y el oxígeno.
\end{solution}

\begin{solution}{exo:g4:air-a-mixture-of-gases:2}
Unas cuatro quintas partes del \omterm{def:g4:air-a-mixture-of-gases:air}{aire} son nitrógeno, y más o menos una
quinta parte es oxígeno.
\end{solution}

\begin{solution}{exo:g4:air-a-mixture-of-gases:3}
Unos 78 litros de nitrógeno y unos 21 litros de oxígeno.
\end{solution}

\begin{solution}{exo:g4:air-a-mixture-of-gases:4}
La botella no está vacía: está llena de \omterm{def:g4:air-a-mixture-of-gases:air}{aire}, y el \omterm{def:g4:air-a-mixture-of-gases:air}{aire} atrapado impide
que entre el agua.
\end{solution}

\begin{solution}{exo:g4:air-a-mixture-of-gases:5}
Hay menos oxígeno (16 partes de cada 100 en lugar de 21) y más dióxido
de carbono (4 partes de cada 100 en lugar de casi nada).
\end{solution}

\begin{solution}{exo:g4:air-a-mixture-of-gases:6}
El oxígeno en los dos casos: tanto una llama como nuestro cuerpo usan
el oxígeno del \omterm{def:g4:air-a-mixture-of-gases:air}{aire}.
\end{solution}

\begin{solution}{exo:g4:air-a-mixture-of-gases:7}
78 cuadrados son nitrógeno; $100 - 78 = 22$ cuadrados no lo son.
\end{solution}

\begin{solution}{exo:g4:air-a-mixture-of-gases:8}
La llama gasta el oxígeno del \omterm{def:g4:air-a-mixture-of-gases:air}{aire} del tarro. Cuando queda demasiado
poco oxígeno alrededor de la llama, se apaga.
\end{solution}

\begin{solution}{exo:g4:air-a-mixture-of-gases:9}
Una quinta parte de 500 litros: $500 \div 5 = 100$ litros de oxígeno.
\end{solution}

\begin{solution}{exo:g4:air-a-mixture-of-gases:10}
Las personas del aula inspiran oxígeno y espiran dióxido de carbono. Al
abrir las ventanas vuelve a entrar \omterm{def:g4:air-a-mixture-of-gases:air}{aire} fresco, con más oxígeno y menos
dióxido de carbono.
\end{solution}

\begin{solution}{exo:g4:air-a-mixture-of-gases:11}
El triple de \omterm{def:g4:air-a-mixture-of-gases:air}{aire} contiene el triple de oxígeno: unos
$3 \times 10 = 30$ segundos.
\end{solution}
